\section{Week 11: Actual Encoder Trace Closure}

Week~11 closes \claimid{OBL-RANS-ENCODE-REFINEMENT} as
\status{proved, success-conditioned partial correctness}.  The decisive
theorem opens the actual generated
\procname{RansEncodeTarget.M._rans_encode}; it is not a wrapper or a pure-only
encoder theorem.  Decoder refinement, the core inverse, and the full HBZ
wrapper remain \status{partial}.

\subsection*{Why the Week 10 local invariant was insufficient}

The Week~10 predicate remembered an arbitrary array snapshot immediately
before one normalization loop.  It could prove a local store and frame, but
lost the normalization bytes written by earlier outer iterations.  The actual
encoder needs the continuation-aware relation
\[
  \mathsf{segment\_matches}\bigl(
    encp,off_0-k,
    \mathsf{innerWritten}(x_0,s,k)
      \mathbin{++}
    \mathsf{encodeTrace}(tail).\mathsf{bytes}
  \bigr).
\]
Week~11 therefore replaces the arbitrary snapshot with predicates relative to
the initial encoder array and the exact pure tail of the actual symbol array.

\subsection*{Tail-aware inner invariant and exact exit}

For outer index $j$, let
\[
  tail=\mathsf{symbolSuffix}(symbols_0,j+1),\qquad
  x_0=\mathsf{encodeTrace}(tail).\mathsf{state}.
\]
The compiled \procname{encoder_inner_tail_inv} carries
\[
  0\le k\le total\le2,
  \quad x=\mathsf{innerStateAfter}(x_0,k),
  \quad off=off_0-k,
\]
the concatenated byte segment above, and
\(
  \mathsf{prefixFrame}(enc_0,encp,off_0-k)
\).
Each actual generated decrement, \texttt{set8}, and right shift preserves this
predicate.  At guard exit,
\procname{encoder_inner_tail_exit_exact} uses the actual
\(x_{\max}\le x\) guard and the mode-2 normalization bound to prove
\[
  k=\mathsf{mode2NormalizationLen}(x_0,s),
\]
so the written list is exactly
\(\mathsf{mode2NormalizationBytes}(x_0,s)\).

\subsection*{Actual table-load and word-step bridge}

The generated weakest-precondition exposes literal loads for \(x_{\max}\),
reciprocal frequency, bias, packed complement/frequency, and shift.  The
theorem family in
\theoryname{Mode2RansEncoderGeneratedWordStep} proves those loads, the
high-half W64 product, the nested shift ladder, and the two W32 additions equal
the certified pure update:
\[
  \mathsf{generatedUpdate}(x,s)
  =\mathsf{W32.ofInt}\bigl(\mathsf{hbzFastEncodeStep}(\mathsf{uint}(x),s)\bigr).
\]
The proof discharges the thirteen concrete symbol cases once in that theory.
Erased declassification and \texttt{protect} calls are handled according to
the extracted semantics; they introduce no new assumption.

\subsection*{One-symbol transition and failure branch}

The compiled success transition combines
\begin{align*}
  \mathsf{encodeTrace}(s::tail).\mathsf{state}
    &=\mathsf{fastStep}(\mathsf{normalize}(x_0,s),s),\\
  \mathsf{encodeTrace}(s::tail).\mathsf{bytes}
    &=\mathsf{normBytes}(x_0,s)
      \mathbin{++}\mathsf{encodeTrace}(tail).\mathsf{bytes}
\end{align*}
with the actual inner exit, word update, cursor arithmetic and prefix frame.
It reconstructs the complete outer live invariant for index $j$.

If the actual body finds $off<4$, it writes $bad=1$.  The loop invariant
then deliberately stops requiring a trace equality.  The final theorem retains
this genuine failure disjunct instead of assuming encoder success.

\subsection*{Whole outer-loop closure}

The live outer invariant is
\begin{align*}
  0 &\le i\le1024,\\
  \mathsf{uint}(x)
    &=\mathsf{encodeTrace}(\mathsf{symbolSuffix}(symbols_0,i)).\mathsf{state},\\
  \mathsf{uint}(off)
    +|\mathsf{encodeTrace}(\mathsf{symbolSuffix}(symbols_0,i)).\mathsf{bytes}|
    &=1024,
\end{align*}
together with exact segment matching, an initial-array prefix frame, and
\(4\le\mathsf{uint}(off)\).  It initializes at $i=1024$ with state
$2^{23}$, empty trace, and cursor 1024.  A live outer-loop exit implies
$i=0$, so the accumulated suffix belongs to
\(
  \mathsf{symbolListOfArray}(symbols_0)
\).

\subsection*{Final serialization prepend}

The generated procedure subtracts four from $off$ and performs four
little-endian byte stores.  The compiled finalization theorem turns their
weakest-precondition expression into
\[
  \mathsf{serialize32LE}(\mathsf{uint}(x))
    \mathbin{++}
  \mathsf{encodeTrace}(\mathsf{symbolListOfArray}(symbols_0)).\mathsf{bytes}
  =
  \mathsf{traceBytes}(\mathsf{symbolListOfArray}(symbols_0)).
\]
It simultaneously preserves the initial prefix below the returned offset.

\subsection*{Compiled generated-procedure theorem}

The exact precondition of
\procname{actual_rans_encode_trace_closure} is
\begin{itemize}
  \item initial encoder, state and symbol arrays are named snapshots;
  \item the concrete mode-2 HBZ encoder table is used;
  \item (count=1024); and
  \item the actual symbol array satisfies the mode-2 canonical stream
        predicate.
\end{itemize}
Its postcondition is
\[
 bad\ne0\;\lor\;
 \left(
 \begin{array}{l}
 bad=0\ \land\ 0\le off\le1020\ \land\ 4\le1024-off\le1024,\\
 \mathsf{segmentMatches}
   (encp,off,\mathsf{traceBytes}(\mathsf{symbolListOfArray}(symbols_0))),\\
 off+|\mathsf{traceBytes}(\mathsf{symbolListOfArray}(symbols_0))|=1024,\\
 \mathsf{prefixFrame}(enc_0,encp,off)
 \end{array}
 \right).
\]
The corollary
\procname{actual_rans_encode_trace_refinement} exposes this full suffix result
without an existential trace and without success in the precondition.  Both
theories fresh-compile with \texttt{-no-eco}.

\subsection*{Rejected tactics and proof decomposition}

A one-shot \texttt{inline; sim} was not retried: reverse outer iteration and a
variable-length inner loop do not preserve the needed pure continuation
automatically.  Expanding the generated reciprocal shift ladder inside the
outer proof created an unmanageable expression; the dedicated generated-word
theorem was rewritten instead.  The initial W64 cursor was normalized under
its canonical range explicitly.  No observation wrapper, procedure copy,
arbitrary trace witness, success premise, authored axiom, or 1024-fold
unrolling was added.

\subsection*{Partial correctness, non-vacuity, and next gate}

The theorem is a Hoare partial-correctness result: every terminating execution
satisfies the displayed failure/success disjunction.  It does not prove local
encoder losslessness, termination, or a terminating all-6 success witness.
Thus \claimid{OBL-RANS-ACTUAL-SUCCESS-WITNESS} stays \status{partial}.

\begin{center}
\begin{tabular}{l l}
\toprule
Claim & Week 11 status \\
\midrule
\claimid{OBL-RANS-ACTUAL-INNER-NORMALIZE} & \status{proved} \\
\claimid{OBL-RANS-ENCODE-REFINEMENT} & \status{proved/partial correctness} \\
\claimid{OBL-RANS-ACTUAL-SUCCESS-WITNESS} & \status{partial} \\
\claimid{OBL-RANS-DECODE-REFINEMENT} & \status{partial} \\
\claimid{OBL-RANS-CORE-INVERSE} & \status{partial} \\
\claimid{OBL-SIG-HBZ-ENCODE-DECODE} & \status{partial} \\
\bottomrule
\end{tabular}
\end{center}

No decoder success, HBZ round trip, canonical parse, encoding delta, or
security result is claimed.  The Week~12 decision is \status{GO-ENCODER}: the
single next obligation is the actual generated decoder trace refinement.
